\chapter{The dataset}
\label{impl:dataset}
\emph{Mirrors Chapter~\ref{ch:dataset} section by section, opening with the data
acquisition (Sec.~\ref{sec:acquisition}) and including the pre-processing and preliminary-
characterization sections (Secs.~\ref{sec:preproc},~\ref{sec:prelim}) -- the cleaning and
first-look steps are part of describing the dataset, not of the planned analysis in
Chapter~\ref{ch:global}.}

\flow{Table~\ref{tab:impl-status} collects the implementation status of every step
discussed below; the sections point to it instead of repeating status remarks.}

\begin{flowblock}
\begin{table}[htbp]
\centering
\footnotesize
\renewcommand{\arraystretch}{1.2}
\begin{tabular}{|l|l|l|}
\hline
\textbf{Pipeline step} & \textbf{Status} & \textbf{Where} \\
\hline
Data acquisition (fetch, download, verify) & implemented & Sec.~\ref{impl:acquisition} \\
\hline
IGC parsing and format validity & implemented, unit-tested & Sec.~\ref{impl:igcformat} \\
\hline
Catalog and season index (\texttt{build-catalog}) & implemented & Sec.~\ref{impl:catalog} \\
\hline
Census scan, diagnostic figures and macros & implemented & Sec.~\ref{impl:preproc}, Table~\ref{tab:impl-repro} \\
\hline
Glider-label normalisation & implemented, unit-tested & Sec.~\ref{impl:glider} \\
\hline
Altitude-channel choice (\texttt{alt\_source}) & \rev{implemented, unit-tested} & Sec.~\ref{impl:altchannel} \\
\hline
Fix-level cleaning (three detectors) & \rev{implemented, unit-tested} & Sec.~\ref{impl:fixlevel} \\
\hline
Validation battery for the cleaning & planned, follows the cleaning & Sec.~\ref{impl:fixlevel} \\
\hline
Trimming of ground phases, interior guard & \rev{implemented, unit-tested} & Sec.~\ref{impl:trimming} \\
\hline
Flight-level filtering & \rev{implemented, unit-tested} & Sec.~\ref{impl:flightfilter} \\
\hline
ENU conversion & \rev{implemented, unit-tested} & Sec.~\ref{impl:enu} \\
\hline
Uniform sampling and segmentation & \rev{implemented, unit-tested} & Sec.~\ref{impl:uniform} \\
\hline
Smoothing and differentiation & \rev{implemented; timescales measured} & Sec.~\ref{impl:savgol} \\
\hline
Preliminary characterization & \rev{implemented} & Sec.~\ref{impl:prelim} \\
\hline
\rev{MSD audit (residuals, controls, isotropy)} & \rev{implemented} & \rev{Sec.~\ref{impl:obs-global}} \\
\hline
\rev{Ensemble MSD (Sec.~\ref{sec:obs-global})} & \rev{implemented, unit-tested} & \rev{Sec.~\ref{impl:obs-global}} \\
\hline
\end{tabular}
\caption{Implementation status of the pipeline steps, collecting the status remarks
of this appendix in one place. \emph{Implemented} code is committed and in use;
\emph{designed} steps have their rules and thresholds fixed (and, where noted, their
parameters already in the configuration) but no running routine yet; \emph{planned}
steps follow the designed ones.}%
\label{tab:impl-status}
\end{table}
\end{flowblock}

\chapter{Data acquisition}
\label{impl:acquisition}
\emph{Mirrors Chapter~\ref{ch:acquisition}.}

\section{Pipeline and tooling}
\label{impl:acqpipeline}
Package \texttt{soaring.acquisition.ffvl}, exposed as two CLI entry points,
\texttt{soaring-para} and \texttt{soaring-delta}, sharing subcommands: \texttt{fetch-xml}
(archive season XMLs), \texttt{download} (fetch IGC files), \texttt{build-catalog}
(rebuild catalog/index), \texttt{status} (per-season coverage), \texttt{verify}
(post-hoc integrity scan, Sec.~\ref{impl:acqintegrity}).

\paragraph{Resumable downloads.} Files already present on disk are skipped, so an
interrupted job resumes freely without re-fetching.

\paragraph{Atomic writes.} Each file is written to a temporary \texttt{.part} and then
renamed, so a file under its final name is always complete -- important on the exFAT
external disk, which has no journaling.

\paragraph{Content validation.} A response is accepted only if it looks like a real IGC
file (an \texttt{A} header record followed by at least one \texttt{B} fix record); HTML
error pages and truncated responses are rejected and retried.

\paragraph{Courteous networking.} A small bounded thread pool, per-request timeouts,
exponential backoff, and a minimum delay between requests avoid hammering the server; a
browser TLS fingerprint is impersonated to pass the site's Cloudflare challenge.

\section{Reproducibility}
\label{impl:acqrepro}
Destination directory (\texttt{data\_root}) set per machine through an environment
variable (\texttt{SOARING\_PARA\_DATA\_ROOT}/\texttt{SOARING\_DELTA\_DATA\_ROOT}) or a
config file shipping a placeholder (\texttt{configs/para\_download.yaml}/
\texttt{configs/delta\_download.yaml}); every command validates the setting at start-up.

\section{Integrity}
\label{impl:acqintegrity}
The \texttt{verify} subcommand re-scans every downloaded file on disk, reusing the same
IGC validation applied at download time, and flags truncated or invalid files, leftover
\texttt{.part} fragments, and read errors.


\section{\rev{Where the data lives}}
\label{impl:onDisk}

\begin{revblock}
Everything is on one external disk and nothing is in the repository, so a reader who has
the disk can locate any number this thesis quotes. There is one root per discipline,
identically laid out, found through an environment variable rather than a hard-coded
path (\path{SOARING_PARA_DATA_ROOT}, \path{SOARING_DELTA_DATA_ROOT}); the repository keeps
only the small versioned \path{data/<discipline>/seasons_index.csv}, so the season tables
of Sec.~\ref{sec:coverage} build without the disk mounted.

\begin{center}
\footnotesize
% sizes are the paraglider root's; the hang-glider one has the same shape, with ~25x
% fewer retained flights and ~37x fewer bytes
\begin{tabular}{|>{\raggedright\arraybackslash}p{4.6cm}|l|p{6.1cm}|}
\hline
\textbf{Path} & \textbf{Size} & \textbf{What it is} \\
\hline
\multicolumn{3}{|l|}{\itshape raw/ --- immutable, never modified after download} \\
\hline
\path{raw/igc/<season>/} & \num{186052} files &
  the track files, one directory per season, named \path{<date>_<flight_id>.igc} \\
\hline
\path{raw/raw_xml/<year>.xml} & 27 files &
  the season listings exactly as the site served them: the provenance of the catalog \\
\hline
\multicolumn{3}{|l|}{\itshape catalog/ --- regenerable from raw/, metadata only} \\
\hline
\path{catalog/catalog.csv} & \SI{87}{\mega\byte} &
  one row per listed flight, 23 columns (Sec.~\ref{sec:catalog}) \\
\hline
\path{catalog/seasons_index.csv} & \SI{2}{\kilo\byte} &
  one row per season: listed, with a track, downloaded \\
\hline
\multicolumn{3}{|l|}{\itshape logs/ --- what the acquisition did} \\
\hline
\path{logs/download.log},\newline\path{logs/failures.csv} & \SI{80}{\kilo\byte} &
  the run log and the retry registry, one row per track that could not be fetched \\
\hline
\multicolumn{3}{|l|}{\itshape derived/ --- regenerable from raw/, everything the analysis makes} \\
\hline
\path{derived/track_scan.parquet} & \SI{8.5}{\mega\byte} &
  the census cache: one scalar row per parsed flight (Sec.~\ref{impl:preproc}) \\
\hline
\path{derived/fixes.parquet} & \SI{43}{\giga\byte} &
  the processed trajectories: $1.36\times10^{9}$ rows for paragliders \\
\hline
\path{derived/segments.parquet} & \SI{4}{\mega\byte} &
  one row per segment the splitting produced, retained or not \\
\hline
\path{derived/flights_meta.parquet} & \SI{10}{\mega\byte} &
  one row per flight \emph{attempted}, dropped ones included \\
\hline
\end{tabular}
\end{center}

The three levels are a rule and not a habit: \path{raw/} is never written to after the
download, \path{catalog/} is a function of \path{raw/}, and \path{derived/} is a function
of both. Deleting \path{derived/} costs computation and never data. It also means the two
census families cannot be confused: \path{track_scan.parquet} describes \emph{parsed}
tracks and is what the \texttt{\textbackslash StatScan*} macros query, because a cut must
be audited on the population it acts on; \path{flights_meta.parquet} describes what the
pipeline \emph{did}, and is what the \texttt{\textbackslash StatPipe*} macros query.

\paragraph{The trajectory table.} \path{fixes.parquet} carries one row per grid point of
every retained segment, keyed \path{(source, flight_id, segment_id)}. Printed as it is on
disk (the hang-glider root; \path{scripts/reporting/show_dataset.py} produces this view):

{\footnotesize
\begin{verbatim}
34,389,245 rows in 43 row groups, 1,156.0 MB on disk, ZSTD

dtypes:
  source  str      t  float32      v_E  float32      interpolated     bool
  flight_id str    E  float32      v_N  float32      edge             bool
  segment_id int16 N  float32      v_z  float32      hampel_flagged   bool
                   z  float32      a_*  float32      alt_invalidated  bool

head(4), columns 1-9:
    source flight_id  segment_id    t          E           N           z       v_E       v_N
hangglider       975           0  0.0  -2.644085   -4.948205 1458.585693 13.275144  1.215108
hangglider       975           0 10.0 111.172691  -33.921207 1460.657104  9.778760 -5.975531
hangglider       975           0 20.0 198.742096 -103.775276 1468.514282  8.025670 -6.961107
hangglider       975           0 30.0 279.220947 -138.911942 1463.914307  7.008325 -2.623670

head(4), columns 10-17:
      v_z       a_E       a_N       a_z  interpolated  edge  hampel_flagged  alt_invalidated
-0.532143 -0.436803 -1.029317  0.192857         False  True           False            False
 0.721429 -0.262474 -0.408811  0.057857         False  True           False            False
 0.625000 -0.088144  0.211696 -0.077143         False False           False            False
-1.466667  0.072790  0.463077 -0.047143         False False            True            False
\end{verbatim}
}

\noindent The four rows are the first four grid points of a \SI{10}{\second} logger:
$t$ starts at zero, $E$ and $N$ within metres of it, $z$ at its measured
\SI{1458}{\metre}. Coordinates are \texttt{float32}, a part in $10^{7}$ on a coordinate
of tens of kilometres and orders below the GPS noise on it. The file is written in
batches of 400 flights, and an analysis streams it rather than loading
\SI{43}{\giga\byte}---which is how the mean-squared displacement of
Sec.~\ref{sec:obs-global} is computed.

\rev{A row group is not a batch, and the difference cost a measurement. The writer hands
Parquet 400 flights at a time; Parquet then cuts what it is handed into row groups of
its own default size, which is why the header above reads \num{43} row groups for an
archive written in \num{16} batches. A reader that takes a row group for a unit of
flights therefore meets, at each boundary, one flight arriving as two fragments---and
processes it as two flights. The ensemble MSD counted \num{157134} paraglider flights
where the pipeline had written \num{156017}, \SI{0.7}{\percent} too many; worse, the
time-averaged MSD is computed \emph{within} a segment and cut at half its length, so a
straddling segment stopped answering at half of its longer fragment instead of half of
itself, and the longest segments---the ones that carry the long-lag tail, and the ones
likeliest to straddle---were the ones truncated. The reading was thus biased by the
writer's buffering rather than by the data. Every consumer now goes through one reader,
\path{soaring.analysis.derived.stream_flights}, which holds back the last flight of each
row group and prepends it to the next, so what it yields is always a whole flight; it
raises rather than continues if the writer's ordering ever stops holding. The dataset
verifier reads through it too, since it had the same blind spot in a worse place: the
step \emph{joining} two fragments was never differenced, so it was never checked, and a
check with a blind spot reports success over it.}

The three timing conventions a reader of this table has to know are argued elsewhere in
the chapter and worth restating together: $t$ is zero at the first fix of free flight
\emph{of the parent flight}, so a segment never restarts it (Sec.~\ref{sec:uniform});
$E$ and $N$ are zero at that same fix (Sec.~\ref{sec:enu}); and $z$ is the adopted
altitude at its measured value, never re-zeroed. The only times in another clock are
\path{ground_phase_start_s} and \path{ground_phase_end_s} in \path{flights_meta}, which
are in the recorded clock of the IGC file because the integrity gate counts its removals
inside that window.

\paragraph{The segment table.} One row per segment the splitting produced, retained or
not, and the reason for each drop:

{\footnotesize
\begin{verbatim}
    source flight_id  segment_id  t_start   t_end  n_fix  n_fix_raw  frac_interpolated
hangglider       975           0      0.0 20640.0   2065       2064           0.001453
hangglider      1032           0      0.0 14175.0      0        691           0.001479

  censored_start  censored_end   kept                 drop_reason
           False         False   True                        <NA>
           False         False  False channel_not_reconstructable
\end{verbatim}
}

\noindent \path{n_fix} counts the rows the segment contributed to \path{fixes} (zero when
dropped) against \path{n_fix_raw}, its measured fixes. \path{segment_id} is assigned at
split time and stays stable, so a gap in the numbering is itself the record of a drop.

\paragraph{The flight table.} One row per flight \emph{attempted}, forty-one columns, read
down rather than across---one retained flight beside one the flight-level filter rejected:

{\footnotesize
\begin{verbatim}
                     a retained flight           a dropped one
source                      hangglider              hangglider
flight_id                          975                     830
pipeline_version                 1.1.0                   1.1.0
drop_stage                         NaN           flight_filter
drop_reason                        NaN  duration_below_minimum
alt_source                        baro                    gnss
n_fix_raw                         2099                     349
n_fix_clean                       2069                     348
n_removed_frozen                    30                       0
n_flagged_kept                      40                       8
n_boundaried                         0                       0
integrity_fraction                 0.0                0.004587
ground_phase_start_s             300.0                   485.0
duration_flight_s              20645.0                  1085.0
path_km                     283.408578                6.044969
lat0                         43.812717                     NaN
dt_native_s                       10.0                     NaN
n_segments_kept                    1.0                     NaN
savgol_window_horiz                5.0                     NaN
\end{verbatim}
}

\noindent (abridged: the table carries \num{47} columns in all.) A dropped flight
keeps everything the stages \emph{before} the verdict had already measured---flight 830
was judged on a duration of \SI{1085}{\second}, and its cleaning counters remain
auditable even though none of its fixes reached the trajectory table. That is what makes
the removals a census rather than an absence.

\paragraph{What is not stored.} Anything that is pure algebra of the stored columns---the
speed $|\mathbf{v}|$, the heading, the angular velocity, the turn radius, the glide ratio,
the mechanical energy---is computed at analysis time and never materialised. Only the
outputs of a non-trivial parametric step are kept, which is what makes the table's size a
function of the archive rather than of the questions asked of it.

\paragraph{Checking it.}\label{impl:verify} \path{scripts/verify_dataset.py} reads the
three tables back and tests, on the data itself, the properties this chapter claims:
completeness within a segment, a strictly increasing and uniform time base, no impossible
step between two interior samples of a segment (Sec.~\ref{sec:fixlevel}), the origin at
$r(0)=0$, reachability---$|\mathbf{r}(t)|\le v_{xy}^{\max}\,t$ for every sample---and
referential integrity between the tables. Each check names the section it comes from, so
a failure says which claim is false rather than only that something is wrong.

Reachability is the one check stated without reference to anything the pipeline did: it
follows from the frame and the speed bound alone, because the flight began at the origin.
That is what makes it able to catch a defect no per-step test can see. A flight whose
\emph{first} fix is corrupt has every step plausible---it simply sits, in the archive,
\SI{4500}{\kilo\metre} from where it says it began, with the whole trajectory referred to
that point (Sec.~\ref{sec:fixlevel}).

The qualification in the third check is not a hedge but the statement being precise.
The bound is a claim about \emph{measured} motion, and two kinds of sample in the table
are not measurements. At a segment edge the smoothing polynomial is evaluated off-centre
(Sec.~\ref{sec:savgol}), so its position carries more variance than a centred one and
on a slow logger can imply a step over the bound where the underlying fixes do not.
Across a bridged gap the monotone cubic is a reconstruction, and its slope may exceed
the secant it spans---a hole crossed at \SI{42}{\metre\per\second}, itself under the
bound and so legitimately bridged, leaves interpolated points implying nearly
\SI{60}{\metre\per\second}. Both kinds are flagged in the table for exactly this
reason (\path{edge}, \path{interpolated}), so the check reads the invariant between
samples that are measured and centred, and \emph{reports} the excess on the others.
That reported number is one of the acceptance criteria Sec.~\ref{sec:savgol} asks for:
measured, rather than assumed away.

The script found two defects this text would otherwise have asserted without warrant.
A randomised robustness suite (\path{test_pipeline_robustness.py}, which combines
injected spikes, null-island runs, frozen locks, time defects, gaps and offsets at
random and asserts only these invariants) found two more, and both were faults in the
\emph{statement} rather than in the pipeline: the invariant had been written without
either qualification above. A third, the corrupt opening block of
Sec.~\ref{sec:fixlevel}, was found by neither: it was found by the ensemble MSD of
Chapter~\ref{ch:global} being wrong, which is the expensive way to find a bug, and
the reachability check above is what makes it cheap next time. It also found one in the
cleaning---a track crossing the
antimeridian is read as motionless, because the great-circle distance is correctly aware
of the wrap and therefore reads the two numerical extremes of a straddling bounding box
as metres apart. No flight in the present archive reaches within ten degrees of that
meridian; the further sources of Sec.~\ref{sec:acquisition} may.
\end{revblock}

\section{The IGC tracklog format}
\label{impl:igcformat}
The parser (\path{soaring.analysis.igc.parse_igc}) checks, at decode time and
independently of any flight-dynamics judgement: that the time is a valid 24-hour
\texttt{hh:mm:ss}; that the arc-minutes field of latitude and longitude is
\texttt{<\,60} (the \path{DDMM.mmm} encoding is undefined otherwise); that the
hemisphere letter is one of the two valid characters; and that the decoded
latitude/longitude fall within $[-90,90]$/$[-180,180]$ degrees. A record failing any of
these is dropped as unparseable, on the same footing as a record with the wrong length or
a non-numeric field. \flow{This format-validity check and the fix-level plausibility checks are distinct stages,
and both are implemented and unit-tested (Sec.~\ref{impl:fixlevel},
Table~\ref{tab:impl-status}).}

\section{Flight metadata and the catalog}
\label{impl:catalog}
Both \path{catalog.csv} and \path{seasons_index.csv} are regenerated on demand
(Sec.~\ref{sec:catalog}) by the acquisition CLI's \texttt{build-catalog} subcommand
(Sec.~\ref{impl:acqpipeline}) -- never hand-edited.

\section{Glider categories}
\label{impl:glider}
\path{aile\_class} is the FFVL's own code (\enquote{C ou 2}, \enquote{non homologuée}),
not the EN/FAI label of Table~\ref{tab:aileclass}, so it is mapped rather than merely
read: \path{soaring.reporting.glider\_class.canonical\_wing\_class} carries the map, one
dictionary per discipline, and raises if a flight's raw code is not a key of it rather
than letting an unrecognised value enter a stratum under its own spelling. The placeholder
class, the literal \enquote{0} in \path{aile\_class}, and a blank entry both come back as
\texttt{NaN} -- unclassified and excluded from a class stratum, not recovered from the
wing-model field, since that would need a wing-model-to-class lookup this repository does
not have. The mapping runs where the class is joined onto a discipline's flights for a
class-stratified figure or macro (Sec.~\ref{sec:glider}), not at the fix-level stage,
since \path{aile\_class} is catalog metadata (Sec.~\ref{impl:catalog}) rather than a
per-fix quantity. Status: Table~\ref{tab:impl-status}.

\begin{flowblock}
\paragraph{Provenance of the EN/FAI class ladders.} The two systems of
Table~\ref{tab:aileclass} are not documented on equal terms. The FAI classes come from the
Sporting Code~\cite{fai_sc7_2024}, which the FAI publishes for free: anyone can download the
current edition and check a definition against it. The EN classes come from a
\emph{European standard}, which is a different kind of document. Standards are drafted by
the European Committee for Standardization (CEN) and then published and \emph{sold} by the
national standards bodies --- AFNOR in France, DIN in Germany, BSI in the United Kingdom ---
as copyrighted documents, priced per copy at a few hundred euros. There is no public
version: unlike a law or a sporting code, the authoritative text of EN~926-2 cannot be
consulted without buying it, and it cannot be quoted here.

The EN entries of Table~\ref{tab:aileclass} are therefore condensed from Table~1 of the
\emph{final draft} of the standard, \mbox{FprEN~926-2:2012}, the version submitted to
formal vote, which the CEN working group circulated publicly~\cite{fpren926_2_2012}. That
draft is the closest thing to the source text that is openly available, and the class
descriptions in it are unlikely to have moved in the published \mbox{EN~926-2:2013}, though
that has not been checked against the published text.
\end{flowblock}

\section{Coverage and statistics}
\label{impl:coverage}
The season tables of Sec.~\ref{sec:coverage} are typeset from the versioned
\path{seasons_index.csv} (Sec.~\ref{impl:catalog}); the trajectory-level census numbers
quoted from Sec.~\ref{sec:preproc} onward come from a different pipeline, the cached
track scan described at the head of Sec.~\ref{impl:preproc}.

\section{Trajectory pre-processing}
\label{impl:preproc}
\emph{Mirrors Sec.~\ref{sec:preproc} and its subsections below.} Parsing: each IGC
file is decoded by the \path{soaring.analysis.igc} parser, which reads the \texttt{B}
records at their fixed character positions (Sec.~\ref{impl:igcformat}) and returns both
altitude channels alongside time and geodetic coordinates. The Welch/PSD
procedure used to compare altitude channels is stated once, in
Sec.~\ref{impl:altchannel}; the smoothing calibration (Sec.~\ref{impl:savgol}) reuses
it, applied to the horizontal components and to the flight's adopted altitude channel.

\paragraph{The census scan and its cache.}
One full parse of the archive backs every trajectory-census number and figure in this
appendix and in Chapter~\ref{ch:dataset}: the \path{load_or_scan_tracks} scan reads
every parsed track (\num{\StatScanParaTracks} paraglider, \num{\StatScanHangTracks}
hang-glider) once and caches a per-flight summary table at
\path{<data_root>/derived/track_scan.parquet}, one cache per discipline.

\begin{revblock}
That cache is worth describing, since every census number in the chapter is a query
against it. It is a table with \emph{one row per parsed flight} and one column per summary
quantity, all produced by \path{track_stats} in
\path{soaring.analysis.census} as it walks a flight's fixes once:
\begin{itemize}[nosep]
  \item \path{duration_s}, \path{n_fix}, \path{dt_s} --- the recorded span, the number of
        \texttt{B} records, and the median inter-fix interval;
  \item \path{path_km}, \path{extent_km} --- the summed great-circle steps and the
        farthest any fix lies from the first;
  \item \path{max_gap_ratio}, \path{missing_fraction} --- the largest inter-fix gap in
        units of \path{dt_s}, and the share of a uniform grid at that interval with no fix;
  \item \path{max_vxy_mps}, \path{max_vz_mps} --- the extreme inter-fix speeds;
  \item \path{baro_alt_min_m}, \path{baro_alt_max_m}, \path{baro_present_frac} --- the
        barometric range and how much of the flight the channel covers.
\end{itemize}
Nothing in it is a trajectory: it is deliberately one scalar row per flight, which is what
makes it small enough to hold in memory for the whole archive and to re-query in seconds.
Its purpose is exactly that---to decouple the cost of \emph{reading} the archive (hours,
once) from the cost of \emph{asking} it a question (instant, as often as a threshold
changes). Every \texttt{\textbackslash StatScan*} macro, and the flight-level and
sampling diagnostics of Figs.~\ref{fig:preproc} and~\ref{fig:gaps}, are functions of these
columns alone.

\rev{\emph{What the census is not.} It is measured on the \textbf{raw} archive, one parse
of the \texttt{B} records with no cleaning, no trimming and no filtering---which is the
only thing it could be, since its purpose is to choose the thresholds that cleaning,
trimming and filtering then apply, and a threshold chosen on the population that survives
it is chosen circularly. So every \texttt{\textbackslash StatScan*} number in this thesis
is a statement about the archive \emph{as recorded}: a duration is the recorded span and
not the airborne one, a path includes whatever a corrupt fix contributed, and a
retained-fraction curve is what a cut would keep if applied alone to the raw population.
The retained ensemble has its own numbers and its own family, \texttt{\textbackslash
StatPipe*}, computed from \path{flights_meta.parquet} after all seven stages have run
(Table~\ref{tab:pipecensus}); where a quantity appears in both, the two are not expected
to agree, and the difference is precisely what the pipeline did. The two families are
never mixed in one statement.}
\end{revblock}

\path{generate_census_stats.py} recomputes the quoted macros from that cache against
the thresholds currently in \path{configs/preprocessing.yaml}, so a threshold change
propagates to every quoted count by re-running one script, with no rescan. \rev{The cache
is refreshed when the archive changes---which for the FFVL collections means once a season
closes, and more substantially when a new source is added (the sailplanes of
Sec.~\ref{sec:acquisition} above all); between refreshes the cached counts can trail the
season tables' download counts, and the two are not expected to agree to the flight.} The same scan feeds
Figures~\ref{fig:preproc}, \ref{fig:gaps} and~\ref{fig:sampling}, regenerated together
by one run of \path{generate_preproc_figure.py} (Table~\ref{tab:impl-repro}); the
fix-level diagnostics of Figure~\ref{fig:fixlevel} are the one exception, drawn from a
separate sampled pass (Sec.~\ref{impl:fixsample}).

\subsection{Choice of altitude channel}
\label{impl:altchannel}

\paragraph{\rev{The \texttt{alt\_source} field.}} \rev{The channel a flight uses is stored as
an \path{alt_source} column (\texttt{baro} or \texttt{gnss})}\flow{ --- two-valued and
categorical, so its storage cost is negligible ---}\rev{ in the processed dataset's
per-flight table, alongside the duration and the raw fix count---the number of \texttt{B}
records the IGC file yields at parse time, before any cleaning, since the channel choice
precedes the cleaning in the pipeline order}\flow{.}

\paragraph{\flow{Welch/PSD procedure and its sample (Figure~\ref{fig:altnoise}a--c).}}
\flow{The sample is a seeded random subsample of \num{3000} flights per
discipline.}\rev{ The size is set by convergence of the ensemble
spectrum, not by the spectral shape: with a few hundred flights the high-frequency part
of the median barometric curve is still visibly ragged (under-averaged), and it smooths
out only with a few thousand flights, so \num{3000} buys a stable band at modest cost.
This is an assertion the reader cannot presently check, and it should not stay one: the
convergence figure---the same median spectrum drawn at a few hundred, a thousand and
\num{3000} flights, on one pair of axes---is to be produced and shown here, so that the
choice of sample size rests on a picture rather than on a claim.}
Panel~(a)/(b) use one
representative flight (\rev{the one with the most fixes among the \num{3000} sampled
flights of that discipline}); its measured offset
(panel~b) is annotated on the panel itself, not restated here.

IGC loggers record one fix per
second nominally, but timestamps are not guaranteed to be perfectly regular (a missed
fix, a logger writing at \SI{1.05}{\second} instead of \SI{1.00}{\second}); Welch's
method (Appendix~\ref{app:psd}) requires a uniformly sampled signal, so the following
runs on every flight before the transform.
\begin{enumerate}[label=(\roman*)]
  \item \textbf{Modal sampling period.} For each flight the median inter-fix interval is
        computed and rounded to the nearest integer second. The value that appears most
        often across all flights in the subsample (the \emph{mode}) is taken as the
        common target $\Delta t$ (in practice $\Delta t=\SI{1}{\second}$ for both
        disciplines, giving $f_s=\SI{1}{\hertz}$ and $f_{\text{Nyq}}=\SI{0.5}{\hertz}$).
  \item \textbf{Flight selection.} Only flights whose own median period lies within
        \SI{0.25}{\second} of the modal $\Delta t$ are retained, so pooled spectra share
        one frequency axis. \rev{Mixing cadences here would not be admissible: a spectrum
        is defined on a frequency axis that runs to $1/(2\Delta t)$, so a \SI{1}{\second}
        and a \SI{9}{\second} logger do not even share a Nyquist frequency, and averaging
        their periodograms per frequency bin would compare different bands. That is the
        whole reason for this cut---and it is a cut on the \emph{diagnostic} sample only.
        The pipeline proper keeps every flight at its native cadence
        (Sec.~\ref{sec:uniform}); it is only this spectral estimate, which pools flights,
        that needs them commensurate.} Flights with fewer than $N_{\text{seg}}=256$ fixes
        ($\sim$\SI{4}{\minute} at $\Delta t=\SI{1}{\second}$) are also excluded
        \flow{(a series shorter than one Welch block of $N_{\text{seg}}=256$ samples
        cannot contribute a periodogram, Appendix~\ref{app:psd})}. Of the \num{3000} sampled flights per
        discipline, this leaves \num{1641} paragliders but \num{727} hang gliders:
        paragliders are predominantly \SI{1}{\second} loggers
        (\StatScanParaDtOneSecPct\,\%, Sec.~\ref{impl:uniform}), so the population mode is
        a good match for most of them, whereas hang gliders split across several native
        rates ($1$, $2$, $5$, $\geq\!8\,\text{s}$, Figure~\ref{fig:sampling}) and only
        their \SI{1}{\second}-logger minority survives this cut.
  \item \textbf{Uniform resampling.} The altitude series $z(t)$ is linearly interpolated
        onto a grid $t_0,t_0+\Delta t,t_0+2\Delta t,\ldots$ spanning the flight.
        \rev{Linear and not cubic, because the resampling must not disturb the band the
        estimate is about to read. Interpolation acts on a spectrum as a filter: linear
        interpolation attenuates high frequencies smoothly and monotonically, whereas a
        higher-order interpolant has a transfer function that can exceed unity near the
        Nyquist frequency, adding power that was never in the signal. Since the quantity
        being measured \emph{is} the flat floor at high frequency, an interpolant that can
        lift it would bias the very number in question; one that only attenuates cannot
        manufacture a floor, and the timestamps are near-uniform to begin with, so the
        attenuation it does apply is slight.}
  \item \textbf{Welch on the resampled series.} \path{scipy.signal.welch} is called
        with $f_s=1/\Delta t$, $N_{\text{seg}}=256$, 50\% overlap, \rev{a Hann window of
        the same $N_{\text{seg}}=256$ samples (the window spans one Welch segment by
        construction---\SI{256}{\second} at the \SI{1}{\second} cadence, giving a
        frequency resolution of about \SI{4}{\milli\hertz})}, and
        linear detrending per segment.
  \item \textbf{Robust ensemble estimator.} The per-flight PSDs are kept individually
        (per discipline and channel) and reduced across flights by the per-frequency
        \emph{median}, with a shaded \numrange{10}{90} percentile band; panel~(c) then
        pools the two disciplines, whose medians coincide. The median, not the
        arithmetic mean, because the per-flight spectra are heavy-tailed: a mean swings
        by one to two orders of magnitude between samples -- for the
        hang-glider barometric floor, from \num{0.8} to \num{93}\,\si{\metre\squared\per\hertz}
        between an $n=\num{114}$ and an $n=\num{727}$ draw -- whereas the median is stable.
        Because $\Delta t$ is the same for every retained flight by construction, the
        figure's Nyquist dotted line is a single fixed value.
\end{enumerate}

\paragraph{Reading the band (panel c).}
\begin{description}
  \item[Why the medians coincide.] For the typical flight both channels are dominated at
        high frequency by the metre resolution to which the IGC format rounds every
        altitude: the white-noise floor of a \SI{1}{\metre} quantiser,
        $\tfrac{1}{12}/f_{\text{Nyq}}\approx\SI{0.17}{\metre\squared\per\hertz}$,\rev{%
        \footnote{\rev{Rounding to a step $q$ adds an error that is well modelled as
        uniform on $[-q/2,q/2]$ and independent between samples, of variance $q^2/12$.
        A white sequence of variance $\sigma^2$ has a flat one-sided power spectral
        density $\sigma^2/f_{\mathrm{Nyq}}$, since integrating it over
        $[0,f_{\mathrm{Nyq}}]$ must return the variance. With $q=\SI{1}{\metre}$ and
        $f_{\mathrm{Nyq}}=\SI{0.5}{\hertz}$ this is
        $(1/12)/0.5\approx\SI{0.17}{\metre\squared\per\hertz}$.}}} sits
        just below where the measured medians level off
        ($\approx\SI{0.3}{\metre\squared\per\hertz}$).
  \item[Where the channels part.] In the tail: at the 90th percentile the
        high-frequency GNSS floor is \num{4}$\times$ the barometric one for paragliders
        and \num{108}$\times$ for hang gliders, and \SIrange{13}{20}{\percent} of flights
        carry a GNSS floor above three times their own barometric one. \rev{What sits in
        that tail is receiver tracking noise scaled by the dilution of precision, since a
        raised \emph{floor} is by definition a white contribution; multipath, being
        correlated over minutes (Sec.~\ref{sec:altchannel}), does not show up there. Why
        the hang-glider fleet is worse in this respect is not settled by these spectra
        alone---older or lower-grade receivers are the obvious candidate---and the figure
        is not claimed to answer it.}
  \item[Consequence.] This heterogeneity, rather than a uniform gap, motivates the
        choice of the barometric channel (Sec.~\ref{sec:altchannel}).
\end{description}

\paragraph{\flow{Census (panel d).}} \flow{Panel~(d) is a full census, not a sample: the
barometric-presence fraction is read off the shared track scan (Sec.~\ref{impl:preproc}),
already paid for by the flight-level diagnostics.}

\begin{revblock}
\enquote{Present} needs a number, and the number is a decision. The scan records
\path{baro_present_frac}, the share of a flight's fixes carrying a non-zero pressure
altitude, and a flight counts as barometric when that share reaches
$\mathtt{BARO\_PRESENT\_MIN}=\SI{95}{\percent}$ (defined once, in
\path{soaring.analysis.altitude_noise}, and imported by both the cleaning and the census
so the three cannot drift). The threshold is high deliberately. Presence is bimodal---the
census below shows the channel is either essentially absent or essentially complete---so
any value in the middle of the range separates the same two populations; but a
mid-scale cut would also admit a flight missing, say, a third of its pressure altitudes,
and such a flight is not one the vertical analysis should treat as barometric: its
$v_z$ would rest on a channel with holes, reconstructed by interpolation over stretches
long enough to erase transitions. Setting the cut at \SI{95}{\percent} sends those
flights to the GNSS fallback, where their altitude is at least uniformly sourced. The
number quoted in the body (Sec.~\ref{sec:altchannel}) is the same one.
\end{revblock}

\rev{Presence is whole-file, checked on the same census:} when the channel is absent it reads exactly zero for
\StatScanParaBaroZeroPct\,\% (paragliders) / \StatScanHangBaroZeroPct\,\% (hang gliders) of
those flights, and when present it covers every fix for \StatScanParaBaroFullPct\,\% /
\StatScanHangBaroFullPct\,\% of them\rev{; the rest miss individual fixes, dropped one by
one (Sec.~\ref{sec:fixlevel}), and no more than \SI{5}{\percent} of the flight by
construction, since below that the channel is not counted as present at all}.

\rev{The bimodality asserted above is what these numbers measure. Where the channel is
absent it is \emph{exactly} zero for \StatScanParaBaroZeroPct\,\% of those flights, and
where it is present it is complete for \StatScanParaBaroFullPct\,\%; the population in
between is thin. Raising the threshold from \num{0.5} to \num{0.95} reclassified
\num{241} paraglider flights and no hang-glider ones---\SI{0.13}{\percent} of the
archive---which is the direct measurement of how thin. Those flights are also where the
worst partial coverage sat: the largest number of missing barometric fixes on a flight
still counted as barometric fell from \num{14878} to
\num{\StatScanParaBaroMissMaxFixes} when the threshold moved, so most of what looked like
a pathological tail was a classification artefact rather than a sensor one.}
\flow{The liveness threshold is \texttt{baro\_min\_range\_m} (key \texttt{alt\_channel};
value in Table~\ref{tab:pipelinemap}); rule and rationale in Sec.~\ref{sec:altchannel}.}

\flow{\paragraph{Why GNSS altitude is vertically noisier.} The asymmetry behind the
vertical dilution of precision is geometric: a receiver sees satellites only above the
horizon, so the ranging directions bracket it in the horizontal plane but all arrive from
the same side vertically, leaving the vertical coordinate the worst-determined one. The GPS
performance standard commits to \SI{8}{\metre} horizontal against \SI{13}{\metre} vertical
error at \SI{95}{\percent} globally, and \SI{15}{\metre} against \SI{33}{\metre} at the
worst site~\cite{gps_sps_ps_2020}. Multipath --- the direct signal arriving together with
replicas delayed by reflection off terrain or the pilot's own equipment --- biases the
range on timescales of minutes rather than adding independent noise from fix to
fix~\cite{garciacrespillo2024}, which is why it raises the low-frequency end of the
spectrum rather than the high-frequency floor read in panel~(c) above.}

\flow{\paragraph{What the liveness test does not catch.} The statistic
\texttt{baro\_min\_range\_m} is a single whole-flight $\max-\min$, so a barometer that
records normally for the first ten minutes and then freezes keeps a range of hundreds of
metres, is judged alive, is adopted, and delivers a vertical velocity of identically zero
over most of the flight --- precisely the outcome the test exists to prevent. Nothing
downstream closes the gap: the frozen-lock rule of Sec.~\ref{sec:fixlevel} needs the
\emph{position} to be frozen, and a cruising wing never trips it; the flight-level
altitude-activity cut consumes the same extremum difference and inherits the same
blindness. A test that did close it would be a \emph{local} one --- a moving-window range,
or a longest-constant-run bound --- and that is a new threshold with its own failure mode
against a wing genuinely holding altitude, so it is not adopted here on the strength of an
argument alone. What the present rule guarantees is exactly this: a channel constant
throughout is refused, one that dies partway is not.}

\subsection{Fix-level cleaning}
\label{impl:fixlevel}

This section gives the working definitions the narrative (Sec.~\ref{sec:fixlevel}) states in
words: the three detectors, their parameters, the order they run in, and how the whole step
is validated. In the order of Table~\ref{tab:cleaning}, the three families map onto a
robust local test (a short window of context), a structural rule (a whole run), and an
absolute bound (no context). Status: Table~\ref{tab:impl-status}.

The detectors run in a fixed order, which the subsections below follow:
\begin{enumerate}[label=(\arabic*), itemsep=0pt, topsep=3pt]
  \item format validity, at decode time (Sec.~\ref{impl:igcformat});
  \item time defects: backward steps, then duplicates;
  \item position outliers (\rev{the impossibility gate; the Hampel flag is recorded, not required---Sec.~\ref{sec:fixlevel}});
  \item frozen-lock runs;
  \item altitude: absolute bounds and altitude outliers.
\end{enumerate}
The order matters because a speed or a residual is only meaningful once the time base is
monotone and free of duplicates. \rev{Merging duplicates is componentwise and unweighted:
fixes sharing one UTC second collapse to a single fix at that second, carrying the mean of
their latitudes and longitudes and the mean of each altitude channel over the members that
have one; the \texttt{V} flag of the merged fix is the logical \emph{or} of the members',
so a degraded member cannot be averaged away.} The position-outlier and frozen-lock passes are one
left-to-right scan that carries the last accepted fix as an anchor: it never re-walks the
whole track, and it absorbs a run by extending the anchor across the flagged block --
tracking the block's running bounding box to test its bounding diameter against $\delta_{xy}$ -- so
the cost is $O(N)$ in the number of fixes.

\subsubsection{\flow{Time-base defects}}
The parser unwraps the UTC-midnight roll-over --- a drop of
$\sim$\SI{86400}{\second} advances the day count --- and \emph{leaves} residual backward
jitter in place (Sec.~\ref{impl:igcformat}). That is deliberate, and it is the one behaviour
the redesign changed: the parser used to flatten the jitter with a running maximum, which is
a repair, and repairing it there destroys the evidence this stage has to act on. A backward
step that is \emph{not} a roll-over is now removed rather than absorbed -- the fix with the corrupt
time is deleted as a node, just as a fix with a corrupt position is. A corrupted timestamp
is recorded as a defect rather than silently repaired. Which fix carries the corrupt time
is decided by minimal removal: delete the smallest set whose removal leaves a strictly
increasing clock (the complement of the longest increasing subsequence), so a single fix
with a forward-jumped clock is itself removed, rather than condemning every later fix
against a running maximum. This requires moving the clamp out of the parser:
\path{parse_igc} keeps the roll-over unwrap (a format concern, needing the raw time of day)
but stops flattening backward jitter with the running maximum -- flattened, the evidence the
cleaning must act on no longer exists downstream.

Fixes sharing one UTC second are
collapsed to a single representative at that second's centroid position. This replaces the
provisional \enquote{keep the first} of the earlier design: keeping the first privileges an
arbitrary member and, for a genuine $\ge\!\SI{2}{\hertz}$ logger, would silently halve the
cadence, whereas the centroid is order-independent and the sub-second resolution it discards
is discarded anyway once the flight is resampled to a common $\Delta t$
(Sec.~\ref{impl:uniform}). One dilution case is accepted: if one member of a duplicate pair
is itself a position spike, the centroid halves its amplitude before the outlier pass (which
runs after the time defects) sees it; a large spike is still flagged at half size, and one
small enough to slip under $k\sigma_k$ after halving is within what the smoother absorbs.

\subsubsection{\flow{Robust local-outlier test and action}}
\flow{The test is Eqs.~\eqref{eq:hampelresid}--\eqref{eq:hampel} of
Sec.~\ref{sec:fixlevel}, evaluated on a \emph{time} window of half-width $w$ (not a
sample count), so it adapts to the native cadence and to irregular sampling; working
defaults in Table~\ref{tab:fixlevel-params}. The window must hold at least a handful of
fixes (order five); a sparser neighbourhood falls back to the absolute bounds
alone.}\footnote{\rev{The residual and run-diameter distances are the
haversine great-circle distance of Eq.~\eqref{eq:haversine}; a vectorised off-the-shelf
implementation exists as \path{sklearn.metrics.pairwise.haversine_distances}, should the
hand-rolled one need cross-checking.}}

Three properties distinguish the test from the bare speed bound:
\begin{itemize}[itemsep=2pt, topsep=3pt]
  \item \emph{The floor $\varepsilon_{\min}$ is not cosmetic.} On a very smooth stretch
        $\sigma_k\to 0$, so without it any sub-metre wiggle would exceed $k\sigma_k$ and
        the test would chase GPS noise where there is no outlier.
  \item \flow{\emph{A burst is one block.} A short burst of consecutive spikes inflates
        the windowed MAD only mildly, so it is still flagged, and it is removed as the
        block whose bracketing neighbours rejoin within scale (the attribution argument
        is the removal-policy paragraph of Sec.~\ref{sec:fixlevel}).}
  \item \emph{It does not clip real manoeuvres.} A genuine sharp turn is drawn by
        consecutive, mutually consistent fixes, each well predicted by its neighbours and
        so of small residual; only an isolated point off the local trend scores large.
        \flow{The one failure mode, the under-sampled tight turn
        (Sec.~\ref{sec:fixlevel}), is measured directly by the injected-defect test
        (check~(1) below).}
\end{itemize}
The flag itself is context-sensitive -- a \SI{30}{\metre\per\second} step out of a
\SI{12}{\metre\per\second} glide is flagged while the same step passes unremarked where it
is ordinary -- but a flag alone does not delete (next paragraph): a
flagged-but-physically-possible fix is \emph{kept}, its flag recorded as a per-flight
anomaly count alongside the integrity-gate counters. \rev{The test runs on the horizontal
position \emph{alone}. An earlier draft of this appendix had it running on the adopted
altitude channel as well, which the three-part argument of Sec.~\ref{sec:fixlevel}
supersedes: on $z$ an absolute bound is sufficient where on $xy$ it is not, a local test
on $z$ would flag the thermal entries the segmentation is built on, and a false vertical
flag is expensive exactly where it is most likely. The vertical is therefore cleaned by
its absolute bounds and by nothing else.}

\paragraph{From flag to action.} The identifier only
\emph{flags}; it never repairs. A flagged horizontal fix is \emph{deleted} as a node -- not
overwritten with the local median -- \flow{with all reconstruction deferred to the single
resampling step (Sec.~\ref{impl:uniform}; the one-policy rationale is
Sec.~\ref{sec:fixlevel})}. This is
the \emph{Hampel identifier}, the detection half of the eponymous filter, not
the median-substituting filter itself: the local median enters only as the reference
for the residual $r_k$, never as an output value.

\flow{\paragraph{Why the flag is not required.} The rule was originally a
conjunction --- flagged \emph{and} impossible --- and running it over the archive exposed a
limit in the \emph{flag}. The Hampel identifier has a breakdown point of \SI{50}{\percent}:
past that, the median moves onto the corruption instead of resisting it. Some loggers write
the null position $(0,0)$ for scattered runs of a few seconds each; where several such runs
fall inside one \SI{\PreprocHampelWindowS}{\second} window they exceed half of it, the
median follows them, the corrupt fixes score a residual of \emph{zero} and go unflagged,
and it is the good fixes around them that are flagged instead. On one such flight the
identifier caught 17 of 99 null fixes, and every survivor left a step of some
\SI{5000}{\kilo\metre} in the record. The rule adopted is therefore the impossibility gate
on its own: a genuine tight circle is still never a candidate, because it flies at an
ordinary speed, and the attribution the flag supplied is instead supplied by the rejoining
test, which names the block whose removal restores the trend. The change bites only where
the identifier is outside its own stated range of validity, which is precisely where the
conjunction failed.}

\begin{description}
  \item[Position.] \rev{Deletion is the impossibility gate: a fix goes when it is
        unreachable from the last accepted fix at the absolute $v_{xy}$ bound \emph{and}
        the removal of the block it belongs to lets the track rejoin within that bound.
        The flag is neither necessary nor sufficient---why it cannot be required is
        Sec.~\ref{sec:fixlevel}---and the genuine under-sampled tight turn is kept because
        it flies at an ordinary speed, not because anything exempts it.}
  \item[\rev{Byte-identical, on the parsed table.}] \rev{The frozen-lock witness of
        Eq.~\eqref{eq:frozenlock} is specified on the raw \texttt{B} records, but the
        parser returns decoded floats and keeps no raw text. No raw text is needed: the
        \path{DDMM.mmm} field is a fixed-width integer count of thousandths of an
        arc-minute, so the decode is injective and two records are byte-identical in
        latitude and longitude \emph{iff} their decoded values compare exactly equal. The
        test is therefore \path{lat[i] == lat[i+1] and lon[i] == lon[i+1]} on the parsed
        columns, with exact equality and no tolerance --- a tolerance would turn the
        conclusive signature into the jittering-freeze case it is meant to exclude.}
  \item[Discontinuity.] When no bounded block restores continuity -- a position
        discontinuity, such as a re-acquisition offset, where every removal still leaves
        an impossible step -- the scan stops deleting and marks a split at the step,
        handled like a long gap (Sec.~\ref{impl:uniform}): both sides are genuine, only
        the transition is unknown.
  \item[\rev{How long a block may run.}] \rev{\enquote{Bounded} needs a number, and it
        is the Hampel half-width $w$: a block that has not rejoined within
        \SI{\PreprocHampelWindowS}{\second} is called a discontinuity instead. The
        value is reused rather than introduced, and it also fixes the scale at which the
        two cases part---an offset smaller than $v_{xy}^{\max}w\approx\SI{900}{\metre}$
        becomes reachable from the anchor before the block expires and is removed as an
        excursion rather than split. That is a real limit of the rule and a bounded one:
        at most one block of genuine fixes, at a discontinuity, against the alternative
        of splitting a flight in two for every isolated spike.}
  \item[\rev{The postcondition, and its counter.}] \rev{After the scan, any step
        between two \emph{surviving} fixes that still breaks the bound is made a segment
        boundary (Sec.~\ref{sec:fixlevel}). The number of boundaries this adds is
        recorded per flight as \path{n_boundaried} in \path{flights_meta}, alongside the
        other cleaning counters: it is the audit of a rule that would otherwise repair
        without saying so, and it is expected to be small---a large value would indict
        the scan rather than the data.}
  \item[Altitude.] \rev{No flag reaches it: the vertical is cleaned by the absolute
        bounds of Table~\ref{tab:cleaning} alone, which mark the value missing and never
        touch the fix\flow{ (a deferral: Sec.~\ref{sec:fixlevel})}. Both bounds act on
        the channel the flight \emph{adopted}, whichever it is. The $|v_z|$ bound was
        placed on barometric distributions, but what it expresses is a climb/sink
        envelope of the aircraft, which does not change with the instrument reading it;
        leaving a GNSS-fallback flight with no vertical cleaner at all --- the
        consequence of dropping the local test --- would be the worse reading. It fires
        more often there, and the rate is recorded per \path{alt_source} so the audit
        can see by how much.}
  \item[\rev{Isolated, not sustained.}] \rev{The $|v_z|$ bound marks an altitude
        missing when the record jumps away from it and back --- both adjoining steps past
        the bound, with opposite signs. This is check~(5) below made operational: a
        \emph{run} of super-threshold steps of one sign is a sustained manoeuvre, a
        spiral dive holding \SIrange{15}{20}{\metre\per\second} of real sink, and it is
        counted and left alone rather than censored.}
  \item[\flow{Why not a Hampel test on $z$ too.}] \flow{The obvious symmetric design ---
        a local-outlier test on the altitude, corroborated by the climb-rate bound, on the
        same footing as the horizontal rule --- is wrong here, for three reasons. First,
        an absolute bound is \emph{sufficient} on $z$ and insufficient on $xy$: the
        vertical envelope is narrow and known (no wing in this dataset sustains
        $|v_z|>\SI{\PreprocMaxVSpeedMps}{\metre\per\second}$, no altitude outside
        $[\PreprocMinAltM,\PreprocMaxAltM]\,\si{\metre}$ is sensor-plausible), whereas a
        \SI{50}{\metre} horizontal spike is an impossible \SI{50}{\metre\per\second} at a
        \SI{1}{\second} cadence but an unremarkable \SI{5}{\metre\per\second} at
        \SI{10}{\second}, so the same defect passes or fails the bound depending on the
        logger and needs a local, scale-free test the vertical case does not. Second, a
        Hampel test on $z$ would flag the physics: the residual from a local median is
        largest where the signal turns fastest, and on the altitude channel that means
        thermal entries and exits, where $v_z$ swings by several \si{\metre\per\second}
        within seconds (Sec.~\ref{sec:uniform}) --- exactly the transitions the phase
        segmentation is built on. A horizontal circle, by contrast, can be told from a
        spike by its speed, a corroborating witness the vertical channel lacks: the
        quantity that would have to arbitrate a suspected altitude spike \emph{is} the
        climb rate, the quantity under suspicion. Third, the costs are not symmetric: a
        false vertical flag is interpolated back cheaply, except across a thermal entry,
        where the interpolation erases the very transition it should preserve, so a
        cheap-looking false positive is expensive exactly where it is most likely.
        Keeping the vertical rule to a bound that only fires on the physically impossible
        avoids that failure mode altogether.}
  \item[\rev{The first fix.}] \rev{An anchor-carrying scan never questions its own first
        anchor, and no rule here supplies one: the discontinuity is left as a
        discontinuity, for the reason Sec.~\ref{sec:fixlevel} gives. The guarantee is
        enforced one stage later, by the reach criterion of
        Sec.~\ref{sec:flightfilter}, and checked on the written table
        (Sec.~\ref{impl:verify}).}
  \item[\flow{Two rules tried, both wrong.}] \flow{An impossible step is a statement
        about a \emph{pair} of fixes and says nothing about which endpoint is wrong; away
        from the ends the bounded-removal test settles it by evidence (which removal
        restores the trend), but at the very first fix there is no trend behind the step
        to ask. Two attribution rules were tried regardless, and both failed, in opposite
        directions. One fired on the first impossible step, which accuses the earlier end
        always: applied to an ordinary position spike three seconds into a record, it
        deleted the good fixes and promoted the spike to origin. The other fired on a
        discontinuity the scan itself declares near the start --- a discontinuity
        survives only where no bounded removal after the step restored the trend --- but
        the split is placed at the first fix of the block that could not be rejoined, and
        that block lies \emph{after} the step whenever the corruption simply outlasts $w$:
        a \SI{50}{\second} corrupt run five seconds in therefore deleted five good fixes,
        kept fifty bad ones, and left the origin \SI{5039}{\kilo\metre} out. What such a
        rule would have bought is the recovery of two flights in \num{156017}, against a
        failure mode demonstrated twice.}
\end{description}
This is the \emph{default to keep} rule of Sec.~\ref{sec:fixlevel}: removal requires
positive, corroborated evidence.

\subsubsection{\flow{Frozen-lock detection}}
\begin{flowblock}
The rule, its witnesses and its margins are Eq.~\eqref{eq:frozenlock} and
Sec.~\ref{sec:fixlevel}; parameter values in Table~\ref{tab:fixlevel-params}. Three
points are implementation-level.
\begin{itemize}[itemsep=2pt, topsep=3pt]
  \item \emph{Diameter, not step.} The run's bounding diameter is the discriminating
        statistic because a step-based test would fire on every \SI{1}{\hertz} climb:
        the fix-to-fix step of a circling wing, $v_{xy}\Delta t\approx\SI{10}{\metre}$,
        is of the order of $\delta_{xy}$ itself.
  \item \emph{Robust witness.} $\Delta z_{\mathrm{baro}}$ is evaluated between medians
        of the run's ends, so an altitude spike inside the run -- not yet cleaned, since
        the altitude pass comes last -- cannot fake a climb.
  \item \emph{Bookkeeping.} On a barometric flight the \texttt{V} flag and the zeroed
        GNSS altitude are recorded as corroborating diagnostics only; they fill the
        witness slot on GNSS-fallback flights alone (Eq.~\eqref{eq:frozenlock}).
  \item \rev{\emph{A declaration must speak for the run.} Where the recorder's
        declarations are the witness, they are required on more than half the run's
        fixes rather than on any one of them. A single \texttt{V} flag inside an
        otherwise healthy stretch is a momentary loss of a satellite, which every flight
        has; it is the sustained declaration that testifies to a frozen receiver.}
\end{itemize}
A run the rule condemns is marked as a data gap; the split it induces is realised at the
uniform-sampling stage (Sec.~\ref{impl:uniform}), on the same footing as a native long
gap.

\flow{\paragraph{The margins, quantified.} A circling wing turns at
$R=v^2/(g\tan\phi)$\footnote{\flow{The coordinated-turn relation: in a steady banked turn
the lift $L$ balances gravity vertically, $L\cos\phi=mg$, while its horizontal component
provides the centripetal force, $L\sin\phi=mv^2/R$; dividing the two gives
$\tan\phi=v^2/(gR)$, i.e.\ $R=v^2/(g\tan\phi)$.}}, so the tightest paraglider circles
($v\approx\SI{10}{\metre\per\second}$ at a bank of $\phi\approx35^\circ$, ordinary values
for tight thermalling) give $2R\approx\SI{30}{\metre}$, and hang gliders, which fly
faster, turn wider ($v\approx\SI{13}{\metre\per\second}$ at the same bank gives
$2R\gtrsim\SI{50}{\metre}$). Against the working value of $\delta_{xy}$, the diameter test
misses by a factor of at least $2$ for the tightest paraglider circles and $\gtrsim3$ for
hang gliders, more for anything wider --- ratios to be rescaled if $\delta_{xy}$ moves.
Barometrically, a climb of order \SI{1}{\metre\per\second} sustained over
$\tau_{\mathrm{freeze}}$ moves the barometer by $\gtrsim\SI{60}{\metre}$ against the
flatness tolerance $\delta_z$, an order of magnitude margin. Both would have to fail well
outside their designed range at once, on the same run, for the conjunction to cut a
genuine climb: on a barometric flight the rule is, in this sense, conservative by
construction, and errs towards keeping.}

\flow{\paragraph{Why the duration condition abstains.} A run that is spatially collapsed
and witnessed but lasts less than $\tau_{\mathrm{freeze}}$ is kept, deliberately: at
durations of a few tens of seconds a collapsed run and a genuine slow, horizontally
localized stretch --- a search, the core of a tight climb --- are not distinguishable by
these signatures, and cutting on them would remove the very events the analysis measures.
The cost is bounded by the threshold itself: a kept short freeze lengthens one waiting
time by at most $\tau_{\mathrm{freeze}}$, small against the tail of $\psi(\tau)$ where the
anomalous exponent is read, whereas a wrongly cut climb removes a waiting time outright.
Short freezes that do get through leave a run of interpolated-looking, near-identical
positions, and the $\psi(\tau)$ sensitivity checks sweep $\tau_{\mathrm{freeze}}$
precisely to bound how much of the tail depends on this choice.}
\end{flowblock}

\subsubsection{\flow{Absolute bounds and gates}}
The six numbers of
Table~\ref{tab:cleaning} are the \texttt{FixLevelThresholds} dataclass, which
\path{load_preproc_config} reads from \path{configs/preprocessing.yaml} under the key
\path{fix_level}; nothing is hard-coded in Python.
They are the context-free floor of the design, and are placed on the real per-fix
distributions (Figure~\ref{fig:fixlevel}) rather than guessed. The horizontal-speed bound is
set where the bin-to-bin ratio of a finely binned (\SI{2.5}{\metre\per\second}) histogram
settles at $\approx 1$, i.e.\ where the decay has actually stopped, not where it first visibly
slows ($\approx\SI{40}{\metre\per\second}$ for hang gliders, versus the adopted
\SI{\PreprocMaxHSpeedHangMps}{\metre\per\second}), since a shallow but real decreasing stretch can still be genuine
rare dynamics. \flow{The Hampel and frozen-lock parameters live under the same key;
Table~\ref{tab:fixlevel-params} collects them.}

\paragraph{Flight-level integrity gate.} The pass records how much it changed each flight (or
segment) -- \path{n_fix_raw}, \path{n_fix_clean} and \path{frac_interpolated} in
\path{flights_meta}. \flow{A flight whose deleted-plus-reconstructed fraction exceeds $f$
(\SI{\PreprocIntegrityMaxPct}{\percent}, Table~\ref{tab:fixlevel-params}) is dropped
whole; the counting window (airborne only) and the rationale are Sec.~\ref{sec:fixlevel},
whose threshold sweep validates the adopted value rather than setting it.} This is the
cleaning counterpart of the sampling-regularity
cut (Sec.~\ref{impl:uniform}).

\begin{table}[htbp]
\centering
\small
\renewcommand{\arraystretch}{1.15}
\begin{tabular}{|l|p{9.4cm}|}
\hline
\textbf{Parameter} & \textbf{Role, and the scale that sets it} \\
\hline
\multicolumn{2}{|l|}{\itshape Robust local-outlier test (Hampel)}\\
\hline
$w$ & half-width of the local time window; a few native intervals \\
\hline
$k$ & flag threshold, in robust standard deviations $\sigma_k$ \\
\hline
$\varepsilon_{\min}$ & absolute residual floor; a few times the
  horizontal GPS noise, so a smooth stretch ($\sigma_k\to0$) is not chased \\
\hline
\multicolumn{2}{|l|}{\itshape Frozen-lock rule (Eq.~\ref{eq:frozenlock})}\\
\hline
$\delta_{xy}$ & run-diameter tolerance; a few times the GPS noise,
  below the $2R\approx\SIrange{30}{100}{\metre}$ of a circling climb \\
\hline
$\delta_z$ & barometric-flatness witness; a few times the
  sensor noise and quantisation, $\ll$ the $\gtrsim\SI{60}{\metre}$ a real climb gains
  over $\tau_{\mathrm{freeze}}$ \\
\hline
$\tau_{\mathrm{freeze}}$ & minimum span of a cut run; two
  thermalling periods, so a slow, tight climb is never mistaken for a freeze \\
\hline
\multicolumn{2}{|l|}{\itshape Flight-level integrity gate}\\
\hline
$f$ & maximum deleted-plus-reconstructed fraction; past it the
  flight is being rebuilt, not cleaned, and is dropped whole \\
\hline
\end{tabular}
\caption{\rev{Fix-level cleaning: what each parameter of the robust and structural
detectors does, and the physical scale that fixes it. The \emph{values} are deliberately
not repeated here---they are in Table~\ref{tab:workingparams}, with the rest of the
pipeline's working values, so that there is one place to read a number and one place to
change it. They are fixed a priori and checked a posteriori on the per-rule removal
fractions (Sec.~\ref{sec:fixlevel}), with the injected-defect calibration as the direct
measurement of the error rates. All live as config keys under \texttt{fix\_level} in
\texttt{configs/preprocessing.yaml}, alongside the absolute bounds of
Table~\ref{tab:cleaning}.}}%
\label{tab:fixlevel-params}
\end{table}

\subsubsection{\flow{Diagnostics and validation}}

\paragraph{Diagnostic sample (Figure~\ref{fig:fixlevel}).}\label{impl:fixsample} Computed by
\path{fix_level_distributions} and rendered by \path{make_fixlevel_diagnostics_figure}
(Table~\ref{tab:impl-repro}). Panels~(a)--(c) pool a seeded random sample of \num{15000}
paraglider flights (\num{126723499} consecutive-fix pairs, $\sim$\SI{73}{\second} to compute)
and the complete \num{\StatScanHangTracks}-flight hang-glider population (a full census
at this size,
\num{37915334} pairs, $\sim$\SI{25}{\second}) -- under two minutes total. The large sample is
deliberate: resolving the sparse artefact tail well enough to tell it from real dynamics
needs far more points than a headline statistic would. The $x$-range in (a)/(b) extends to
each discipline's own 99.99th percentile rather than cropping just past the bound; panel~(c)
keeps a narrower range, its cut sitting far out in an otherwise tight, unimodal distribution.
Panel~(b) bins one integer \si{\metre\per\second} per bin because barometric altitude is
logged to the metre: differencing at a \SI{1}{\second} cadence returns integer rates, and a
finer grid would alias them into a spurious comb.\footnote{Not every flight samples at
\SI{1}{\second} (Figure~\ref{fig:sampling}), and differencing over a longer native step
returns finer values (multiples of \SI{0.2}{\metre\per\second} at \SI{5}{\second}) that
partly fill the gaps once cadences are pooled: among \SI{1}{\second}-native fixes about
99\%\ land within \SI{0.05}{\metre\per\second} of an integer, against only 34--42\%\ at every
other cadence.}

\paragraph{Validating the cleaning.} \flow{The primary validation is the two-stage
procedure of Sec.~\ref{sec:fixlevel}: freeze the working values a priori, then audit the
per-rule removal fractions over the archive and sweep each threshold around its working
value. The checks below are the fuller battery behind it, run as the implementation
matures.}
\begin{enumerate}[label=(\arabic*), itemsep=2pt, topsep=3pt]
  \item \flow{\emph{Injected defects.} The synthetic-defect procedure of
        Sec.~\ref{sec:fixlevel}, item~(d), gives the detection floor -- the smallest
        spike caught -- and, where the removal-fraction audit flags a rule, recalibrates
        $k$, $w$, $\delta_{xy}$ and $\tau_{\mathrm{freeze}}$ to a target rate: the only
        direct measurement of the cleaning's false negatives.}
  \item \emph{Downstream invariance.} Recompute the transport estimators -- the
        displacement-scaling exponent and the waiting-time-tail exponent -- over a grid of
        cleaning strengths\rev{, i.e.\ re-running the cleaning with all thresholds
        tightened or loosened together by a common factor around the working point,} and take the operating point on the plateau where they stop
        moving; movement under stricter cleaning is the signature of a cut biting into
        signal.
  \item \emph{Detector nesting.} The fixes flagged by the absolute bound should nest
        inside those flagged by the local test (the impossible is a subset of the locally
        anomalous); a large disagreement points to a defect class neither was designed for
        -- a cheap guard.
  \item \emph{Per-phase removal rates.} The error this pipeline most wants to avoid is
        losing a slow-phase fix. Once the segmentation is in place
        (Ch.~\ref{ch:global}), the removal and split rates are measured \emph{per
        phase} and required to be flat across transition, search and climb rather than
        enriched in the two slow phases; and synthetic tight climbs -- circles at flying
        speed, over a range of radii and wind drift -- planted in otherwise clean tracks
        confirm directly that the frozen-lock detector never splits a real climb. This is
        the false-\emph{positive} complement of check~(1), and its natural readout is the
        tail of $\psi(\tau_w)$: a cut biting into climb or search would thin that tail
        before it moved anything else.
  \item \emph{Isolation of the vertical cuts.} The fixes cut by
        $|v_z|>\SI{\PreprocMaxVSpeedMps}{\metre\per\second}$ must be \emph{isolated}, not consecutive: a
        sustained spiral dive holds \SIrange{15}{20}{\metre\per\second} of real sink, so a
        run-shaped excess is the signature of a genuine manoeuvre -- calling for a raised
        bound or a coherent-run exemption rather than censoring -- whereas an isolated
        super-threshold difference, bracketed by ordinary rates, is a sensor error and
        safe to cut.
\end{enumerate}

\subsection{Trimming of ground phases}
\label{impl:trimming}
\flow{The two thresholds of Eq.~\eqref{eq:trimming} ($v_0$, $T_0$) are the
\texttt{TrimmingThresholds} dataclass, loaded from \path{configs/preprocessing.yaml}
(key \texttt{trimming}); values in Table~\ref{tab:pipelinemap}.}

The interior-ground guard of Sec.~\ref{sec:trimming} keeps
\texttt{interior\_ground\_s} ($T_{\mathrm{ground}}$, \SI{600}{\second}) and
\texttt{ground\_flatness\_m} (\SI{\PreprocGroundFlatnessM}{\metre}) in the same group.
\rev{The flatness is tested on the \emph{detrended} altitude, as Sec.~\ref{sec:trimming}
requires: fit a straight line to the adopted channel over the candidate stint by least
squares, subtract it, and compare the residual's central span---its 5th to 95th
percentile---against \texttt{ground\_flatness\_m}. The linear term is the barometric
drift, which over a stint of $T_{\mathrm{ground}}$ is already of the order of the
tolerance; the residual is what is left of genuine vertical motion, and the central span
rather than the full range because a range grows with the length of the stint and a
central span does not (Sec.~\ref{sec:trimming}). No new config key is needed.}

\begin{revblock}
Building it made one thing explicit that the rule as stated leaves open. Detrending cannot
by itself separate a pressure drift from a steady climb, because both are linear and the
fit removes either completely: a wing thermalling at \SI{1}{\metre\per\second} in still
air holds a ground speed near zero, gains \SI{600}{\metre} over a $T_{\mathrm{ground}}$
stint, and leaves a detrended residual of \emph{zero} --- so the flatness condition alone
would excise it as a mid-flight landing, which is precisely the failure this guard exists
to avoid. What separates the two is the \emph{rate}: drift runs at tens of metres per
hour, a soaring climb two orders of magnitude faster. The fitted slope is therefore
bounded as well, at $\delta_z/\tau_{\mathrm{freeze}}$ --- the frozen-lock rule's own
statement of what counts as motionless, read as a rate, so the guard borrows a number the
configuration already fixes instead of introducing one.
\end{revblock} \flow{Below
$T_{\mathrm{ground}}$ the same predicate is stored as a list of suspect intervals per
flight, which the segmentation later maps onto the phases they touch (consumption of the
flag: Sec.~\ref{sec:trimming}).}

\flow{\paragraph{Why the central span, not the full range.} A full range is a maximum
statistic: over zero-mean noise it grows without bound with the number of samples, as
$2\sigma\sqrt{2\ln n}$. At the metre-scale barometric noise of this archive it reads
\SI{4.6}{\metre} over a 60-sample stint and \SI{7.1}{\metre} over \num{3000}, so against a
\SI{5}{\metre} tolerance the same motionless pilot would pass the test on a one-minute
stint and fail it on a ten-minute one, on noise alone --- and since an interior stint must
last $T_{\mathrm{ground}}$ to be considered at all, it would fail on every stint the rule
is ever offered. The central span carries the same units and reading and sits at
$3.3\sigma$ whatever the stint's length, so $\delta_z^{\mathrm{gr}}$ means the same thing
on a stint of any duration; a genuine climb of hundreds of metres exceeds it either way.}

\flow{\paragraph{The GNSS-fallback abstention.} The intended rule widens the tolerance on
a GNSS-fallback flight to match that channel's noise; the widening is not implemented,
because the factor it needs is a measured ratio between the two channels' short-timescale
noise, and a first measurement on the archive did not reproduce the ordering
Fig.~\ref{fig:altnoise} reports. Until that is settled, the single barometric tolerance is
applied to both channels, which on a noisier one simply means the flatness condition is
not met and the stint is kept --- the conservative half of the intended rule, since the
error this guard exists to avoid is excising a genuine climb. The cost is that a
mid-flight landing on a GNSS-fallback flight is neither excised nor flagged.}

\subsection{Flight-level filtering}
\label{impl:flightfilter}
\flow{The three cuts of Sec.~\ref{sec:flightfilter} are the
\texttt{FlightLevelThresholds} dataclass (key \path{flight_level} in
\path{configs/preprocessing.yaml}); values in Table~\ref{tab:pipelinemap}. The census
measurements behind the path floor and the upper duration bound, and the totals of
Figure~\ref{fig:preproc}, are quoted through macros computed from the shared track scan
(Sec.~\ref{impl:preproc}); the headline duration and path medians are quoted in
Sec.~\ref{sec:prelim}.}

\rev{The plausibility and reach bounds of Sec.~\ref{sec:flightfilter} are not further
\texttt{flight\_level} keys: both read \path{max_horizontal_speed_mps} from the
\texttt{fix\_level} block, so the same number governs a single step, a whole-flight
average and a whole-flight displacement, and the three cannot drift apart. Both are
evaluated on the trimmed track, after the per-block sums of the duration and the path, and
they are the last of the whole-flight cuts because each is a statement about a record the
earlier ones have already accepted. The extent they need is returned by
\path{flight_quantities} alongside the duration and the path, from the same haversine
call, so it costs nothing: the maximum of a vector the function already forms. Unlike the
duration and the path it is \emph{not} summed per block---a displacement from a fixed
origin has nothing to sum---which is why a split does not weaken it.}

\begin{revblock}
The chain behind every number quoted in Sec.~\ref{sec:flightfilter} is short enough to
name in full, so that a count can be traced to the code that produced it:
\begin{itemize}[nosep]
  \item \path{soaring.analysis.igc.parse_igc} decodes one IGC file into a table of fixes;
  \item \path{soaring.analysis.census.track_stats} reduces those fixes to the one
        summary row described in Sec.~\ref{impl:preproc}, and
        \path{scan_tracks} / \path{load_or_scan_tracks} do that for every flight of a
        discipline and cache the result;
  \item \path{fraction_retained} and \path{retention_curve} in the same module evaluate a
        single cut against a column of that table---the first at one threshold, the second
        along a grid of them, which is what the bottom row of Fig.~\ref{fig:preproc} plots;
  \item \path{scripts/reporting/generate_census_stats.py} calls the same functions with
        the thresholds read from \path{configs/preprocessing.yaml} and writes the
        \texttt{\textbackslash StatScan*} macros the text quotes.
\end{itemize}
Two consequences follow, and both matter for reading the numbers. The counts are
\emph{marginal}: each is the effect of one criterion applied alone to the whole parsed
population, not of the cascade. And they are pre-cleaning: the scan runs on parsed tracks,
so a count in Sec.~\ref{sec:flightfilter} is what the criterion removes from the parsed
archive and not from the cleaned one. That is deliberate and is argued for above
(\emph{What the census is not}): a threshold chosen on the population the cleaning has
already reshaped would be chosen against its own effect. The post-cleaning counts are the
cascade of Sec.~\ref{sec:pipelinemap}, and the two answer different questions.
\end{revblock}

\subsection{From geographic to local Cartesian coordinates}
\label{impl:enu}
No implementation content: this step is pure coordinate geometry with no configurable
parameters, and both of its figures (Figure~\ref{fig:ellipsoid}, Figure~\ref{fig:enu-rotation})
are hand-drawn \textsc{TikZ} diagrams, not data plots. \rev{The mathematics behind
Eqs.~\eqref{eq:ecef} and the ENU rotation is derived in Appendix~\ref{app:geodesy}.}

\subsection{Uniform sampling within a flight}
\label{impl:uniform}
The sampling thresholds (key \texttt{sampling}, the \texttt{SamplingThresholds}
dataclass\flow{; values in Table~\ref{tab:pipelinemap}}) act as follows.
\begin{description}
  \item[Split bound.] $g_{\max}=\min\big(\texttt{max\_gap\_factor}\times\Delta t,\,
        \max(\texttt{max\_gap\_seconds},\,2\,\Delta t)\big)$, pinned by the constraint
        chain of Sec.~\ref{sec:uniform}. Its audit is a per-flight count -- the census
        scan records \path{max_gap_ratio} and \path{dt_s}, so the absolute gap is their
        product -- and joins Figure~\ref{fig:gaps} when regenerated.
  \item[Missing fraction.] \flow{Acts per \emph{segment} (rule:
        Sec.~\ref{sec:uniform}); the per-flight value in the census scan is the
        pre-split diagnostic, not the operational gate.}
  \item[Segment gate.] \texttt{min\_segment\_duration\_s} (\SI{\PreprocMinSegmentS}{\second}). It only
        needs to support the smoothing window and the within-segment statistics: a few
        smoothing timescales, never less than $w\,\Delta t$\flow{; the only gate applied
        at segment level (Sec.~\ref{sec:uniform})}.
  \item[Measured vs interpolated.] A grid point counts as measured when a fix lies
        within half a native step of it, interpolated otherwise; the same tolerance
        decides whether a flight is uniform as recorded (no interpolated points) or was
        resampled (\path{was_resampled}).
\end{description} The bound is on the \emph{time base} and is the same for
every channel: an order-of-magnitude smaller $v_z$ buys no wider gap for the altitude, since
the interpolation error scales with curvature, not speed, and a wide altitude gap would erase
the climb--glide transitions that $v_z$ is there to resolve (Sec.~\ref{sec:uniform}).

\paragraph{Segment handling.}
\begin{description}
  \item[Origin and clock.] \flow{A split assigns a \path{segment_id} and nothing else
        (semantics: Sec.~\ref{sec:uniform}).}
  \item[Tables.] One row per fix in \path{fixes}, keyed
        (\path{flight_id}, \path{segment_id}); per-segment records (start/end, fix
        count, \path{frac_interpolated}, censoring flags) in a \path{segments} table
        beside \path{flights_meta}. No separate files; the ensemble MSD at lag $t$ reads
        $\mathbf{r}(t)$ wherever some segment covers $t$.
  \item[Censoring.] Any phase duration truncated at a segment boundary is flagged
        censored, for the $\psi(\tau)$ fits to exclude.
\end{description}

\paragraph{Per-channel fill.} \flow{Across a short gap the horizontal position is
interpolated by monotone piecewise cubics (\path{scipy.interpolate.PchipInterpolator})
and the adopted altitude channel linearly (its within-phase smoothness makes the choice
among low-order methods immaterial), realising the completeness invariant of
Secs.~\ref{sec:fixlevel} and~\ref{sec:uniform}. Each filled grid point is flagged in an
\path{interpolated} boolean column of \path{fixes}, and \path{frac_interpolated} records
the per-segment fraction for the integrity gate.}

\flow{\paragraph{Why linear for the altitude, cubic for the horizontal.} The asymmetry
follows from what each channel is used for. The altitude is differentiated once more than
the horizontal coordinates before it is read --- $v_z$ is the segmentation's discriminant
--- so an interpolation artefact in $z$ propagates directly into a phase assignment,
whereas the same artefact in $E$ or $N$ is absorbed by the smoothing. Linear interpolation
cannot overshoot, and its error over a hole of width $g$ is bounded by
$\max|\ddot z|\,g^2/8$\footnote{\flow{On an interval of width $g$ the difference between a
$C^2$ function $z$ and the chord through its endpoints is
$z(t)-p(t)=\tfrac12\,\ddot z(\xi)\,(t-t_0)(t-t_0-g)$ for some $\xi$ in the interval; the
quadratic factor is largest at the midpoint, where it equals $g^2/4$, giving
$|z-p|\le\max|\ddot z|\,g^{2}/8$.}}, whereas a cubic is more accurate where the second
derivative is smooth and worse exactly where it is not --- the transitions a phase
assignment turns on.}

\flow{\paragraph{The test flight that forced the \texttt{z\_reconstructed} flag.} $g_{\max}$
bounds the interval between two \emph{fixes}, not an interval over which a fix exists but
carries no altitude value; that case is not rare, since the vertical channel can be absent
while the horizontal record is perfect (a logger writing no barometric value, or one the
fix-level cleaning removed). No time gap opens there, so no split is declared, and the
linear fill spans the hole whatever its length. On a test flight missing \SI{500}{\second}
of altitude the filled $z$ departed from the truth by up to \SI{540}{\metre}, the smoothing
differentiated the straight line into a vertical velocity that was pure invention, and the
record read $\texttt{frac\_interpolated}=0$ with \texttt{was\_resampled} false --- the
flight declaring itself uniform as recorded, because the \texttt{interpolated} flag is a
statement about the time base alone. The vertical therefore carries its own per-fix flag,
\texttt{z\_reconstructed}, with the per-segment share \texttt{frac\_z\_reconstructed} and,
per flight, the longest unbroken run \texttt{z\_gap\_max\_s}, reported next to the fraction
because a fraction alone cannot separate a hundred two-second holes, which a straight line
bridges well, from one twelve-minute hole, which it does not. Any analysis of the vertical
excludes on that flag; the horizontal analyses are untouched by it, which is why the flag
is per channel rather than one for all three.}

\flow{Figures~\ref{fig:gaps} and~\ref{fig:sampling} are rendered by
\path{make_gap_diagnostics_figure} and \path{make_sampling_figure} from the shared track
scan (Sec.~\ref{impl:preproc}).}

\paragraph{Gap-ratio discreteness (Figure~\ref{fig:gaps}).} Checked directly: only
\StatScanParaGapIntPct\,\%/\StatScanHangGapIntPct\,\% (para/hang) of largest-gap-ratio
values are exact integers, the rest a
genuinely continuous scatter; the integer peaks visible for hang gliders sit around
$k=3,6,10,15$. Panels (a)/(b) of the figure therefore use linear axes: a logarithmic
axis would compress precisely this small-integer structure, on which the adopted cut
sits.

\paragraph{Native-rate discreteness (Figure~\ref{fig:sampling}).} Checked directly:
$\Delta t$ lands on an exact whole second for \StatScanParaDtIntPct\,\% of paragliders
and \StatScanHangDtIntPct\,\% of hang
gliders. \flow{Per-rate breakdown: paragliders sit mostly
(\StatScanParaDtOneSecPct\,\%) at \SI{1}{\second}; hang gliders spread across the
rates below.}
\begin{flowblock}
\begin{center}
\small
\begin{tabular}{|l|c|c|c|c|c|}
\hline
Native rate & \SI{1}{\second} & \SI{2}{\second} & \SI{5}{\second} & \SI{10}{\second}
  & other \\
\hline
Hang gliders & \StatScanHangDtOneSecPct\,\% & \StatScanHangDtTwoSecPct\,\% &
  \StatScanHangDtFiveSecPct\,\% & \StatScanHangDtTenSecPct\,\% & remainder \\
\hline
\end{tabular}
\end{center}
\end{flowblock}

\subsection{Smoothing and differentiation}
\label{impl:savgol}
\flow{Three decisions are pinned for the implementation (status:
Table~\ref{tab:impl-status}):}
\begin{flowblock}
\begin{itemize}[itemsep=2pt, topsep=3pt]
  \item the filter runs per segment (never across a boundary) with
        \path{scipy.signal.savgol_filter} in \texttt{mode='interp'}, so the terminal
        half-windows are fitted on the edge window rather than padded, and edge samples
        inherit the segment-boundary caution of Sec.~\ref{impl:uniform};
  \item the window is floored at $w=5$, the smallest a cubic fit accepts
        (slow-logger consequences: Sec.~\ref{sec:savgol}, item~iii);
  \item the polynomial order and the three smoothing timescales are the
        \texttt{SavgolParams} dataclass in \path{configs/preprocessing.yaml} (key
        \texttt{savgol}): \path{tau_c_horizontal_s} for $(E,N)$,
        \path{tau_c_vertical_baro_s} and \path{tau_c_vertical_gnss_s} for the vertical
        channel (the flight's adopted altitude, Sec.~\ref{sec:altchannel});
        values in Table~\ref{tab:pipelinemap}.
\end{itemize}
\end{flowblock}

\paragraph{The measured timescales.} The knees were measured with
\path{estimate_savgol_timescales.py}, which applies the Welch procedure of
Sec.~\ref{impl:altchannel} to $E$, $N$ (local equirectangular metres) and to the altitude
channel each flight actually uses, on a seeded sample (908 horizontal flight-components,
719 barometric and 189 GNSS-fallback flights at \SI{1}{\hertz}). The knee rule: the floor
is the median PSD over \SIrange{0.35}{0.5}{\hertz}; $f_c$ is the lowest frequency whose
median PSD falls below three times it. \flow{The measured floors sit at $\approx0.60$
and $\approx\SI{0.26}{\metre\squared\per\hertz}$ against the theoretical quantisation
floors $q^2/12/f_{\mathrm{Nyq}}$ of $0.57$ and $0.17$; the three knees coincide at
$f_c\approx\SI{0.2}{\hertz}$, hence $\tau_c=\SI{5}{\second}$ for every channel
(Sec.~\ref{sec:savgol}, item~(iv), for the reading), and even the 90th-percentile
spectra move the knee by less than a second.} The two vertical keys stay, conditioned on
\path{alt_source}, for the
noisy-minority refinement; the validation of Sec.~\ref{sec:savgol} (flat residual
spectrum, stability under one window step) is what confirms the values. The filter runs
on the resampled grid, so the small interpolated fraction (\path{frac_interpolated}) is
already in place when it reads its window.

\section{Preliminary characterization}
\label{impl:prelim}
\begin{revblock}
\path{scripts/reporting/generate_prelim_figure.py} produces the figures of
Sec.~\ref{sec:prelim} and of Sec.~\ref{sec:strata-compat} (Chapter~\ref{ch:global}) and the
\texttt{\textbackslash StatPrelim*} macros. It reads three
things: the per-flight summary and the per-lag positions \path{audit_msd.py} wrote
(Sec.~\ref{impl:obs-global}), \path{flights_meta.parquet} for the take-off coordinates,
and the season catalogue for \path{wing_class} and \path{season}. The join on
\path{flight_id} resolves for \SI{99.99}{\percent} of retained flights.

Reading the audit arrays rather than the fix table is what makes the stratification
affordable: a stratified MSD is then a row selection on a matrix already in memory, so the
three stratifications and the isotropy check together cost one pass that has already been
paid, instead of one \SI{43}{\giga\byte} scan each.
\end{revblock}

\paragraph{\rev{The basemap.}} \rev{The take-off panels are drawn on Natural Earth
admin-0 country polygons (public domain), which give the coastline and the borders in one
layer and carry the overseas departments inside the France feature---which is what puts
La~R\'eunion on the map at all. The geometry is not fetched at draw time and Cartopy is
not a dependency: \path{scripts/reporting/build_basemap.py} downloads it once, keeps the
rings that intersect each panel, decimates them (\SI{0.01}{\degree} for France,
\SI{0.001}{\degree} for La~R\'eunion, \SI{0.25}{\degree} for the world frame) and writes
\path{data/basemap.json}, which is committed at \SI{0.76}{\mega\byte}. The figure then
needs nothing but Matplotlib and regenerates from a clean checkout with no network. Rings
are kept whole rather than clipped to the frame---Matplotlib clips to the axes anyway, and
clipping a polygon properly would need the geometry library the arrangement exists to
avoid. Re-run the builder only to change a panel's extent.}

\paragraph{\rev{The three frames.}} \rev{Metropolitan France
$\lambda\in[-5.5,10]$, $\varphi\in[41,51.5]$; La~R\'eunion $\lambda\in[55.15,55.92]$,
$\varphi\in[-21.42,-20.82]$; the world panel for everything outside both. The split is
the archive's, not a choice: \num{\StatPrelimParaReunionFlights} paraglider flights launch
on La~R\'eunion and \num{\StatPrelimParaElsewhereFlights} outside France altogether, and a
single frame that held them would render metropolitan France four pixels wide. Each panel
sets its aspect to $1/\cos\varphi$ at its own mid-latitude, so the coastline has its true
shape rather than a stretched one.}

\paragraph{\rev{The orographic boxes.}} \rev{The strata are longitude/latitude boxes
rather than polygons traced from the map: a box that is written down can be checked
against a gazetteer, and the test the strata support is whether the curves separate, not
where exactly the Alps end. Alps $\lambda\in[5.4,10.0]$, $\varphi\in[43.8,46.6]$; Pyrenees
$\lambda\in[-1.9,3.3]$, $\varphi\in[42.0,43.5]$; Massif Central $\lambda\in[1.8,4.6]$,
$\varphi\in[43.6,46.2]$; anything else inside $\lambda\in[-6,10]$,
$\varphi\in[41,52]$ is \enquote{lowland} and anything outside it is \enquote{abroad}. The
boxes are coarse and overlap along their edges, with the first match winning in the order
listed; they are a stratification, not a geography, and the test they support is whether
the curves separate rather than where exactly the Alps end.}

\paragraph{\rev{The stratum floor.}} \rev{A stratum below \num{2000} flights is pooled
into the residual group rather than drawn. Below that the curve's own sampling error at
the long-lag end is comparable to the separations being tested for, and a panel that draws
it invites the reader to compare noise. The floor is why the hang-glider archive, at
\num{\StatPrelimHangFlights} flights, contributes only to the isotropy panel: one wing
class and one orographic group clear it, so the stratified panels would compare a stratum
against itself.}

\paragraph{\rev{The isotropy test.}} \rev{The variance ratio is computed on the same
per-lag position matrix, column by column, over the flights covering that lag. The
Kolmogorov--Smirnov comparison runs at four lags spanning the fit range, on the signed
components rather than their magnitudes, so that a difference in \emph{spread} and a
difference in \emph{shape} both register. Only the distance is quoted. At
\num{\StatPrelimParaFlights} flights the two-sample $p$-value underflows to zero for any
departure worth naming, so it separates nothing; the distance is the effect size and is
what the text reads.}

\paragraph{\rev{Stratum departure.}} \rev{Each stratum's curve is divided by the pooled
curve over the fit range and the largest $|{\cdot}-1|$ recorded. Two summaries are
generated: the maximum over strata (\texttt{\textbackslash StatPrelim*MaxDeparturePct})
and the median over strata (\texttt{*TypicalDeparturePct}). The maximum alone reads as a
statement about the smallest and noisiest stratum, which is what it usually is; the median
is what distinguishes an axis along which the ensemble genuinely splits from one along
which a single small group wanders. Both are quoted in Sec.~\ref{sec:prelim}.}
